% ============================================================
% Section 4: Design
% ============================================================
% Target length: ~1.5 columns (800–1000 words)
% Write: Month 5 (R1 + R2)
% This is the core technical contribution section
% Figures needed: architecture diagram, token format,
%                 data flow diagram
% ============================================================

\section{Design}
\label{sec:design}

\subsection{Design Principles}

\ferrite{} is built around three principles:

\begin{enumerate}
  \item \textbf{Single privileged boundary.} The capability
    broker is the sole component with authority to grant access
    to sensitive resources. No ambient authority exists anywhere
    in the system.

  \item \textbf{No trust in untrusted code.} Web content,
    extensions, and agent outputs are treated as potentially
    hostile. Capability requests from these components are
    evaluated against policy before any action is taken.

  \item \textbf{Complete auditability.} Every capability grant,
    denial, and exercise is recorded in a tamper-evident log.
    The log provides cryptographic proof that every agent action
    was (or was not) authorized.
\end{enumerate}

\subsection{Architecture Overview}

\todo{Insert architecture diagram here.
      Figure should show: Trusted Zone (broker, policy engine,
      audit log, UI shell) vs Untrusted Zone (Servo, extension
      sandbox, agent runtime, network stack).
      Show capability request flow across the boundary.}

\begin{figure}[h]
  \centering
  % \includegraphics[width=\columnwidth]{figures/architecture.pdf}
  \todo{Figure: architecture diagram — generate from the
        architecture doc diagrams}
  \caption{\ferrite{} architecture. The capability broker is
           the sole privileged boundary between trusted and
           untrusted components.}
  \label{fig:architecture}
\end{figure}

\subsection{Capability Token System}

A capability token is an unforgeable, scoped, time-bounded
authorization credential.
Formally, a token $T$ is a tuple:

\[
T = \langle id, P, C, O, R, t_{exp}, \sigma \rangle
\]

where:
\begin{itemize}
  \item $id$ is a globally unique token identifier (UUID v7)
  \item $P$ is the principal identity (extension, agent, or web content)
  \item $C \in \mathcal{C}$ is the capability type from the
    capability lattice $\mathcal{C}$
  \item $O$ is the origin scope (URL prefix or wildcard)
  \item $R \in \{\text{Low}, \text{Medium}, \text{High}\}$
    is the risk classification
  \item $t_{exp}$ is the absolute expiry timestamp
  \item $\sigma$ is the broker's Ed25519 signature over the
    token fields
\end{itemize}

A token is valid iff: (1) $\sigma$ verifies under the broker's
public key, (2) $t_{now} < t_{exp}$, (3) the request URL
matches $O$, and (4) the policy engine evaluates the token
as allowed.

\subsection{Capability Types}
\label{subsec:capability_types}

\todo{Table of Wave 1 and Wave 2 capability types.
      Copy from CLAUDE.md capability table.}

\subsection{Policy Engine}

\ferrite{} uses Regorus~\cite{regorus}, a pure-Rust implementation
of the Open Policy Agent (OPA) Rego language, as its policy engine.
Policies are composable Rego rules loaded from a directory hierarchy
at broker startup, with a deny-by-default base policy and per-principal
override files.

\todo{Show example Rego policy snippet. Include the default.rego
      deny-by-default rule from policies/default.rego.}

\subsection{Risk Classification and Step-Up Consent}

Capabilities are classified into three risk tiers:
Low (read-only, auto-approved), Medium (write or sensitive read,
logged and auto-approved), and High (credentials, clipboard,
downloads -- requires explicit user consent).

High-risk requests cause the broker to return a
\texttt{ConsentRequired} response rather than a grant.
The UI shell renders a consent dialog in the trusted browser chrome,
isolated from web content.
The agent receives no token until the user explicitly approves.

\subsection{Verifiable Audit Log}

\todo{Describe hash chain structure and Merkle tree upgrade.
      Include the hash computation formula.
      Reference Certificate Transparency lineage.}

\subsection{Extension Sandbox}

\todo{Describe Extism/Wasmtime sandbox.
      Explain capability-gated host functions.
      Explain why ungranted host functions are not linked.}

\subsection{Agent Protocol}

\todo{Describe JSON-RPC over WebSocket agent protocol.
      This section written in Month 4 after agent runtime is built.}
